% GENERATED FILE. Edit canonical lessons, then run npm run generate:books.
% canonical-source-hash: fnv1a64:082ea9a550d1adf2
% canonical-lessons: GU-C226-dhrujvu, GU-C226-haphvu, GU-C226-sughvu, GU-C226-galvu, GU-C226-phephsu

\chapter{Breath nose and throat}
\label{ch:gu-breath-nose-and-throat}
\hlchaptermodality{226}

\begin{chapteropening}
\textbf{By the end of this chapter:} \emph{I can say that I am shivering or out of breath and ask about swallowing.}

Complete the last lesson of chapter 226: I can say that I am shivering or out of breath and ask about swallowing.
\end{chapteropening}

\section[\texorpdfstring{\gu{ધ્રૂજવું} (dhrūjvũ)}{ધ્રૂજવું (dhrūjvũ)}]{\gu{ધ્રૂજવું} (dhrūjvũ) — to shiver, to tremble}

\label{lesson:GU-C226-dhrujvu}



\begin{quote}
Before the new one: say the Gujarati for to yawn, then the Gujarati for sweat.
\end{quote}

\subsection*{You'll want to know: \gu{ધ્રૂજવું}}
\textbf{\gu{ધ્રૂજવું}} — \emph{dhrūjvũ} — \textquotedblleft{}to shiver, to tremble\textquotedblright{}.

\begin{quote}
\textbf{\gu{શું} \gu{તમે} \gu{ઠંડીથી} \gu{ધ્રૂજો} \gu{છો}?} — \emph{śũ tame ṭhanḍīthī dhrūjo cho?} — \textquotedblleft{}Are you shivering with cold?\textquotedblright{}
\end{quote}

\begin{cousinweb}
\textbf{\gu{ધ્રૂજવું}} (\emph{dhrūjvũ}) is probably inherited. CDIAL lists a Gujarati \emph{dhujvũ}, \textquotedblleft{}to tremble, be shaken\textquotedblright{}, under Sanskrit \emph{dhūyate}, \textquotedblleft{}is shaken\textquotedblright{} (CDIAL 6863); the \emph{r} of \emph{dhrūjvũ} is not explained there, so the link is uncertain.
\end{cousinweb}

\begin{grammarlens}[title={What you've built}]
One more word for everyday life at A2.
\end{grammarlens}

\subsection*{Your turn}
\begin{itemize}
\raggedright
  \item \emph{Say it:} \emph{dhrūjvũ}
  \item \emph{Say it:} \emph{dhrūjvũ}, once more
  \item \emph{Say it:} \emph{śũ tame ṭhanḍīthī dhrūjo cho?} — Are you shivering with cold?
  \item \emph{Recall:} say the Gujarati for to yawn, then the Gujarati for sweat, then say \emph{dhrūjvũ} again
\end{itemize}

\begin{tcolorbox}[breakable,colback=teal!4,colframe=teal!35!black,title={Before you move on}]
What does \gu{ધ્રૂજવું} mean? (\textquotedblleft{}To shiver, to tremble\textquotedblright{}.) Say it once more.
\end{tcolorbox}
\section[\texorpdfstring{\gu{હાંફવું} (hā̃phvũ)}{હાંફવું (hā̃phvũ)}]{\gu{હાંફવું} (hā̃phvũ) — to pant, to be out of breath}

\label{lesson:GU-C226-haphvu}



\begin{quote}
Before the new one: say the Gujarati for sweat, then the Gujarati for to shiver.
\end{quote}

\subsection*{You'll want to know: \gu{હાંફવું}}
\textbf{\gu{હાંફવું}} — \emph{hā̃phvũ} — \textquotedblleft{}to pant, to be out of breath\textquotedblright{}.

\begin{quote}
\textbf{\gu{દાદા} \gu{હવે} \gu{નથી} \gu{હાંફતા}.} — \emph{dādā have nathī hā̃phtā} — \textquotedblleft{}Grandfather is not out of breath now.\textquotedblright{}
\end{quote}

\begin{cousinweb}
\textbf{\gu{હાંફવું}} (\emph{hā̃phvũ}) comes from \emph{*hamph}, \textquotedblleft{}pant\textquotedblright{}, an imitative reconstructed form; CDIAL lists the Gujarati verb, glossed \textquotedblleft{}to pant\textquotedblright{} (CDIAL 13973).
\end{cousinweb}

\begin{grammarlens}[title={What you've built}]
One more word for everyday life at A2.
\end{grammarlens}

\subsection*{Your turn}
\begin{itemize}
\raggedright
  \item \emph{Say it:} \emph{hā̃phvũ}
  \item \emph{Say it:} \emph{hā̃phvũ}, once more
  \item \emph{Say it:} \emph{dādā have nathī hā̃phtā} — Grandfather is not out of breath now.
  \item \emph{Recall:} say the Gujarati for sweat, then the Gujarati for to shiver, then say \emph{hā̃phvũ} again
\end{itemize}

\begin{tcolorbox}[breakable,colback=teal!4,colframe=teal!35!black,title={Before you move on}]
What does \gu{હાંફવું} mean? (\textquotedblleft{}To pant, to be out of breath\textquotedblright{}.) Say it once more.
\end{tcolorbox}
\section[\texorpdfstring{\gu{સૂંઘવું} (sū̃ghvũ)}{સૂંઘવું (sū̃ghvũ)}]{\gu{સૂંઘવું} (sū̃ghvũ) — to sniff, to smell}

\label{lesson:GU-C226-sughvu}



\begin{quote}
Before the new one: say the Gujarati for to shiver, then the Gujarati for to pant.
\end{quote}

\subsection*{You'll want to know: \gu{સૂંઘવું}}
\textbf{\gu{સૂંઘવું}} — \emph{sū̃ghvũ} — \textquotedblleft{}to sniff, to smell\textquotedblright{}.

\begin{quote}
\textbf{\gu{શું} \gu{કૂતરો} \gu{ફૂલ} \gu{સૂંઘે} \gu{છે}?} — \emph{śũ kūtro phūl sū̃ghe che?} — \textquotedblleft{}Is the dog sniffing the flower?\textquotedblright{}
\end{quote}

\begin{cousinweb}
\textbf{\gu{સૂંઘવું}} (\emph{sū̃ghvũ}) is a tadbhava. CDIAL lists it, glossed \textquotedblleft{}smells\textquotedblright{}, under Sanskrit \emph{śiṅghati}, \textquotedblleft{}smells\textquotedblright{}, and its variant \emph{*śṛṅkhati} (CDIAL 12579).
\end{cousinweb}

\begin{grammarlens}[title={What you've built}]
One more word for everyday life at A2.
\end{grammarlens}

\subsection*{Your turn}
\begin{itemize}
\raggedright
  \item \emph{Say it:} \emph{sū̃ghvũ}
  \item \emph{Say it:} \emph{sū̃ghvũ}, once more
  \item \emph{Say it:} \emph{śũ kūtro phūl sū̃ghe che?} — Is the dog sniffing the flower?
  \item \emph{Recall:} say the Gujarati for to shiver, then the Gujarati for to pant, then say \emph{sū̃ghvũ} again
\end{itemize}

\begin{tcolorbox}[breakable,colback=teal!4,colframe=teal!35!black,title={Before you move on}]
What does \gu{સૂંઘવું} mean? (\textquotedblleft{}To sniff, to smell\textquotedblright{}.) Say it once more.
\end{tcolorbox}
\section[\texorpdfstring{\gu{ગળવું} (gaḷvũ)}{ગળવું (gaḷvũ)}]{\gu{ગળવું} (gaḷvũ) — to swallow}

\label{lesson:GU-C226-galvu}



\begin{quote}
Before the new one: say the Gujarati for to pant, then the Gujarati for to sniff.
\end{quote}

\subsection*{You'll want to know: \gu{ગળવું}}
\textbf{\gu{ગળવું}} — \emph{gaḷvũ} — \textquotedblleft{}to swallow\textquotedblright{}.

\begin{quote}
\textbf{\gu{શું} \gu{બાળક} \gu{દવા} \gu{ગળે} \gu{છે}?} — \emph{śũ bāḷak davā gaḷe che?} — \textquotedblleft{}Does the child swallow the medicine?\textquotedblright{}
\end{quote}

\begin{cousinweb}
\textbf{\gu{ગળવું}} (\emph{gaḷvũ}) is a tadbhava, a word inherited from Sanskrit and reshaped by sound change. It comes from \emph{galati}, \textquotedblleft{}swallows\textquotedblright{}, beside \emph{gilati} (CDIAL 4075).
\end{cousinweb}

\begin{grammarlens}[title={What you've built}]
One more word for everyday life at A2.
\end{grammarlens}

\subsection*{Your turn}
\begin{itemize}
\raggedright
  \item \emph{Say it:} \emph{gaḷvũ}
  \item \emph{Say it:} \emph{gaḷvũ}, once more
  \item \emph{Say it:} \emph{śũ bāḷak davā gaḷe che?} — Does the child swallow the medicine?
  \item \emph{Recall:} say the Gujarati for to pant, then the Gujarati for to sniff, then say \emph{gaḷvũ} again
\end{itemize}

\begin{tcolorbox}[breakable,colback=teal!4,colframe=teal!35!black,title={Before you move on}]
What does \gu{ગળવું} mean? (\textquotedblleft{}To swallow\textquotedblright{}.) Say it once more.
\end{tcolorbox}
\section[\texorpdfstring{\gu{ફેફસું} (phephsũ)}{ફેફસું (phephsũ)}]{\gu{ફેફસું} (phephsũ) — lung}

\label{lesson:GU-C226-phephsu}



\begin{quote}
Before the new one: say the Gujarati for to sniff, then the Gujarati for to swallow.
\end{quote}

\subsection*{You'll want to know: \gu{ફેફસું}}
\textbf{\gu{ફેફસું}} — \emph{phephsũ} — \textquotedblleft{}lung\textquotedblright{}.

\begin{quote}
\textbf{\gu{શું} \gu{તેનું} \gu{ફેફસું} \gu{નબળું} \gu{છે}?} — \emph{śũ tenũ phephsũ nabaḷũ che?} — \textquotedblleft{}Is his lung weak?\textquotedblright{}
\end{quote}

\begin{cousinweb}
\textbf{\gu{ફેફસું}} (\emph{phephsũ}) is a tadbhava, a word inherited from Sanskrit and reshaped by sound change. It comes from \emph{phupphusa}, \textquotedblleft{}lungs\textquotedblright{} (CDIAL 9090).
\end{cousinweb}

\begin{grammarlens}[title={What you've built}]
One more word for everyday life at A2.
\end{grammarlens}

\subsection*{Your turn}
\begin{itemize}
\raggedright
  \item \emph{Say it:} \emph{phephsũ}
  \item \emph{Say it:} \emph{phephsũ}, once more
  \item \emph{Say it:} \emph{śũ tenũ phephsũ nabaḷũ che?} — Is his lung weak?
  \item \emph{Recall:} say the Gujarati for to sniff, then the Gujarati for to swallow, then say \emph{phephsũ} again
\end{itemize}

\begin{tcolorbox}[breakable,colback=teal!4,colframe=teal!35!black,title={Before you move on}]
What does \gu{ફેફસું} mean? (\textquotedblleft{}Lung\textquotedblright{}.) Say it once more.
\end{tcolorbox}
